\documentclass[11pt, a4paper]{article}

% --- PACCHETTI ---
\usepackage[utf8]{inputenc}
\usepackage[italian]{babel}
\usepackage[margin=2.5cm]{geometry} % Imposta i margini della pagina
\usepackage{graphicx}               % Per inserire il logo
\usepackage{tabularx}               % Per tabelle larghe quanto la pagina
\usepackage[table]{xcolor}          % Per colorare le celle delle tabelle
\usepackage[hidelinks]{hyperref}    % Per stile link cliccabili
\usepackage{lmodern}      % Rende i font scalabili all'infinito
\usepackage{anyfontsize}  % Permette di usare dimensioni personalizzate

\begin{document}

% ==============================================================================
% FRONTESPIZIO (PAGINA 1)
% ==============================================================================
\begin{titlepage}
    \centering
    \vspace*{0.5cm}

    % LOGHI
    \includegraphics[width=0.45\textwidth]{../../../../assets/logo_unipd.jpg}
    \hfill
    \includegraphics[width=0.45\textwidth]{../../../../assets/logo.png}\par
    \vspace{0.5cm}

    % NOME DEL GRUPPO
    {\fontsize{45}{50}\sffamily\selectfont \textbf{On-By-One}\par} 
    \vspace{1cm}

    % TITOLO DEL VERBALE
    {\huge \bfseries Verbale interno 2026/10/23 \par}
    
    \vspace{1.5cm}

    % MEMBRI DEL GRUPPO
    {\large
    \textbf{On-By-One}:\par
    \vspace{0.3cm}
    \begin{tabular}{l l}
        Sara Ballardin & 2101048 \\
        Shuhui Chen & 2138001 \\
        Simone Giuliano & 2145591 \\
        Alvise Greggio & 2141004 \\
        Daniele Nuzzaci & 2139039 \\
        Federico Temporin & 2146764 \\
        Tommaso Zorzetto & 2148629 \\
    \end{tabular}
    \par}

    \vfill

    % TABELLA INFORMAZIONI DOCUMENTO E DESTINATARI
    \begin{table}[h]
        \centering
        \arrayrulecolor{gray!50}
        \renewcommand{\arraystretch}{1.5}
        \begin{tabularx}{\textwidth}{|c|c|c|X|}
            \hline
            \rowcolor{gray!10} \multicolumn{4}{|c|}{\textbf{Informazioni documento}} \\
            \hline
            \rowcolor{gray!10} \textbf{Versione} & \textbf{Data} & \textbf{Stato} & \multicolumn{1}{c|}{\textbf{Destinatari}} \\
            \hline
            1.0.0 & 2026/10/27 & Approvato & 
            On-By-One \newline 
            Prof. Tullio Vardanega \newline 
            Prof. Riccardo Cardin \\
            \hline
        \end{tabularx}
        \arrayrulecolor{black}
    \end{table}
    
    \vspace{0.5cm}
    
    % CONTATTI
    {\large Contatti: \href{mailto:onbyone.swe@gmail.com}{onbyone.swe@gmail.com} }
    
\end{titlepage}

\newpage

% ==============================================================================
% REGISTRO DELLE MODIFICHE (PAGINA 2)
% ==============================================================================
\section*{Registro delle modifiche}

\vspace{1cm}
\begin{center}
    \arrayrulecolor{gray!50}
    \renewcommand{\arraystretch}{1.5}
    \begin{tabularx}{\textwidth}{|c|c|c|X|c|}
        \hline
        \rowcolor{gray!10} 
        \multicolumn{1}{|c|}{\textbf{Versione}} & 
        \multicolumn{1}{c|}{\textbf{Data}} & 
        \multicolumn{1}{c|}{\textbf{Autore}} & 
        \multicolumn{1}{c|}{\textbf{Descrizione}} & 
        \multicolumn{1}{c|}{\textbf{Verificatore}} \\
        \hline
        1.0.0 & 2026/10/27 & Sara Ballardin & Aggiunti ID issue mancanti & Simone Giuliano \\
        \hline
        0.1.1 & 2026/10/25 & Sara Ballardin & Stesura del verbale & - \\
        \hline
        0.1.0 & 2026/10/24 & Sara Ballardin & Bozza iniziale & - \\
        \hline
    \end{tabularx}
    \arrayrulecolor{black}
\end{center}
\vfill
\newpage

% ==============================================================================
% INDICE (PAGINA 3)
% ==============================================================================
\tableofcontents
\newpage

% ==============================================================================
% INFORMAZIONI GENERALI
% ==============================================================================
\section{Informazioni generali}

\begin{table}[h]
    \centering
    \arrayrulecolor{gray!50}
    \renewcommand{\arraystretch}{1.5}
    \begin{tabularx}{\textwidth}{|c|>{\centering\arraybackslash}X|}
        \hline
        \rowcolor{gray!10} \multicolumn{2}{|c|}{\textbf{Dettagli della riunione}} \\
        \hline
        \textbf{Data} & 2026/10/23 \\ \hline
        \textbf{Ora inizio} & 15:00 \\ \hline
        \textbf{Ora fine} & 18:30 \\ \hline
        \textbf{Modalità} & Discord \\ \hline
    \end{tabularx}
    \arrayrulecolor{black}
\end{table}

\begin{table}[h]
    \centering
    \arrayrulecolor{gray!50}
    \renewcommand{\arraystretch}{1.5}
    \begin{tabularx}{\textwidth}{|>{\centering\arraybackslash}X|>{\centering\arraybackslash}X|>{\centering\arraybackslash}X|}
        \hline
        \rowcolor{gray!10} \multicolumn{3}{|c|}{\textbf{Partecipanti}} \\
        \hline
        \rowcolor{gray!10} \textbf{Nome e cognome} & \textbf{Ruolo} & \textbf{Presenza} \\
        \hline
        Sara Ballardin & - & Presente \\ \hline
        Shuhui Chen & - & Presente \\ \hline
        Simone Giuliano & - & Presente \\ \hline
        Alvise Greggio & Responsabile & Presente \\ \hline
        Daniele Nuzzaci & - & Presente \\ \hline
        Federico Temporin & - & Presente \\ \hline
        Tommaso Zorzetto & - & Presente \\ \hline
    \end{tabularx}
    \arrayrulecolor{black}
\end{table}

% ==============================================================================
% CORPO DEL VERBALE
% ==============================================================================

\section{Ordine del giorno}
\begin{itemize}
    \item Organizzazione del \textit{repository} e creazione della bacheca \textit{Kanban}.
    \item Definizione delle responsabilità per i ruoli di progetto.
    \item Stesura del preventivo.
    \item Suddivisione dei compiti per la scrittura delle valutazioni aziendali.
\end{itemize}

\section{Discussioni}
La riunione si è aperta definendo l'organizzazione del repository. Il team ha creato una \textit{project board} in stile Kanban per ottimizzare il tracciamento delle attività, procedendo con la creazione e all'assegnazione delle prime issue su GitHub. \newline
Il gruppo ha definito in modo chiaro le responsabilità e i compiti per ciascun ruolo di progetto (Responsabile, Amministratore, Analista, Progettista, Programmatore, Verificatore), garantendo una chiara separazione per le fasi successive. \newline
La strategia per la stesura del preventivo si basa sulla stima del carico di lavoro atteso per ciascun ruolo nelle diverse fasi di sviluppo. Le 20 settimane di progetto sono state suddivise in cinque periodi per ottenere un calcolo più realistico. Il preventivo finale e le tabelle di ripartizione saranno riportati nel documento di Dichiarazione degli Impegni. \newline
Terminata l'analisi dei capitolati, il team si è diviso il lavoro di stesura. I resoconti degli incontri aziendali e le motivazioni delle scelte saranno raccolti in un documento riassuntivo.

\section{Decisioni}
\begin{itemize}
    \item \textbf{VI-2026-10-23.1} Organizzazione del repository e creazione della bacheca Kanban.
    \item \textbf{VI-2026-10-23.2} Definizione delle responsabilità per ogni ruolo di progetto.
    \item \textbf{VI-2026-10-23.3} Ripartizione delle ore di lavoro in base ai ruoli e alle fasi di sviluppo.
    \item \textbf{VI-2026-10-23.4} Assegnazione della stesura dei testi per il documento riassuntivo sui capitolati.
\end{itemize}

\section{To-do}
\begin{center}
    \arrayrulecolor{gray!50}
    \renewcommand{\arraystretch}{1.5}
    \begin{tabularx}{\textwidth}{|c|X|c|c|}
        \hline
        \rowcolor{gray!10} 
        \multicolumn{1}{|c|}{\textbf{ID issue}} & 
        \multicolumn{1}{c|}{\textbf{Task}} & 
        \multicolumn{1}{c|}{\textbf{Assegnatario}} & 
        \multicolumn{1}{c|}{\textbf{Scadenza}} \\
        \hline
        \href{https://github.com/On-By-One-Team/On-By-One.github.io/issues/16}{\#16} & Redazione del documento di Dichiarazione degli Impegni & - & 2026/10/29 \\
        \hline
        \href{https://github.com/On-By-One-Team/On-By-One.github.io/issues/24}{\#24} & Stesura delle valutazioni sui capitolati & Tutto il gruppo & 2026/10/27 \\
        \hline
    \end{tabularx}
    \arrayrulecolor{black}
\end{center}

\end{document}
